\section{Shortest Routes II}
\label{sec:cses-shortest-routes-ii}

\problemstatement{Given an undirected weighted graph of $n$ cities and $m$
roads, answer $q$ queries, each asking the shortest distance between two
cities. Output $-1$ if unreachable. Here $n$ is small (up to a few
hundred).}

\begin{patternbox}
Many shortest-distance queries between arbitrary pairs, with small $n$, is
\textbf{Floyd--Warshall}: compute all-pairs shortest distances once in
$O(n^3)$, then answer each query in $O(1)$.
\end{patternbox}

\subsection*{Relax through every intermediate}

Let $\textit{dist}[i][j]$ start as the direct edge (or $\infty$, and $0$ on
the diagonal). For each candidate intermediate node $k$, update every pair:
$\textit{dist}[i][j]=\min(\textit{dist}[i][j],
\textit{dist}[i][k]+\textit{dist}[k][j])$. After considering all $k$, every
entry is the true shortest distance, because any shortest path's highest-
indexed intermediate is handled in the round for that $k$.

\begin{ideabox}[Allow One More Intermediate at a Time]
The $k$-loop is outermost: after iteration $k$, $\textit{dist}[i][j]$ is the
shortest path using intermediates only from $\{1,\ldots,k\}$. Expanding the
allowed set one node at a time builds all-pairs distances without missing
any path.
\end{ideabox}

\subsection*{Algorithm}

\begin{enumerate}
  \item Initialise the distance matrix from edges (keep the minimum for
        multi-edges; $0$ on the diagonal).
  \item Triple loop $k,i,j$ relaxing through $k$ (guard against $\infty$
        overflow).
  \item Answer each query from the matrix; $-1$ if still $\infty$.
\end{enumerate}

\subsection*{Dry run}

\dryrun{Triangle $1\!-\!2(4)$, $2\!-\!3(1)$, $1\!-\!3(10)$.}

\begin{itemize}
  \item Through $k=2$: $\textit{dist}[1][3]=\min(10,4+1)=5$.
  \item Query $(1,3)\Rightarrow5$.
\end{itemize}

\subsection*{C++ implementation}

\begin{lstlisting}
#include <bits/stdc++.h>
using namespace std;
const long long INF = 1e18;

int main() {
    int n, m, q;
    cin >> n >> m >> q;
    vector<vector<long long>> dist(n + 1, vector<long long>(n + 1, INF));
    for (int i = 1; i <= n; i++) dist[i][i] = 0;
    for (int i = 0; i < m; i++) {
        int a, b; long long w; cin >> a >> b >> w;
        dist[a][b] = min(dist[a][b], w);
        dist[b][a] = min(dist[b][a], w);
    }

    for (int k = 1; k <= n; k++)
        for (int i = 1; i <= n; i++)
            for (int j = 1; j <= n; j++)
                if (dist[i][k] < INF && dist[k][j] < INF)
                    dist[i][j] = min(dist[i][j], dist[i][k] + dist[k][j]);

    while (q--) {
        int a, b; cin >> a >> b;
        cout << (dist[a][b] >= INF ? -1 : dist[a][b]) << "\n";
    }
    return 0;
}
\end{lstlisting}

\begin{complexitybox}
\textbf{Time Complexity.} $O(n^3+q)$. \textbf{Space Complexity.} $O(n^2)$
for the distance matrix.
\end{complexitybox}

\begin{takeawaybox}
All-pairs shortest distance with small $n$ and many queries is
Floyd--Warshall: one $O(n^3)$ triple loop, $O(1)$ queries. The
$k$-outermost order -- allowing one more intermediate each round -- is the
detail that makes it correct.
\end{takeawaybox}
